We now give a verifier class that is easy to learn in the realizable case, but not when there are misspecifications.

\begin{theorem}
There is a verifier class $\V$ with the following two properties. In the realizable case, cooperative debaters can use a randomized strategy with $\ex[\Mi(T)]<9$ for every $T$, every $v^*\in\V$, and every sequence of inputs. In the unrealizable case, for every $k\ge1$ and every $T\ge 3k$, the environment can choose the inputs and draw $v^*\in\V$ and the misspecified rounds at random before the interaction, so that the expected number of misspecified rounds is $k$ and every strategy of the debaters has $\ex[\Mi(T)]\ge\frac32k$.
\end{theorem}

\begin{proof}
We will consider debates of length one. Let $\X=\naturals$, $\C=\{0,1\}$, and let $\ct_t$ and $\cf_t$ be the certificates at step $t$. For $b\in\{0,1\}$, let $h_b$ be the constant hypothesis $h_b(x)=b$, and let $\H=\{h_0,h_1\}$. For every bit $b$ and every function $e:\X\to(\C\times\C)\cup\{\bot\}$, define
\[
v_{b,e}(x,\ct,\cf)=
\begin{cases}
1-b, & (\ct,\cf)=e(x),\\
b, & \text{otherwise},
\end{cases}
\]
and let $\V=\{v_{b,e}:b\in\{0,1\},\ e:\X\to(\C\times\C)\cup\{\bot\}\}$. For each $x$, think of $v(x,\cdot,\cdot)$ as a $2\times2$ table with rows indexed by the \dt's certificate $\ct$ and columns by the \df's certificate $\cf$. The table of $v_{b,e}$ at $x$ has every entry equal to $b$, except the entry $e(x)$, which equals $1-b$. If $e(x)=\bot$, there is no such entry, and the whole table is $b$. Now $v_{b,e}$ validates $h_b$. The exceptional entry, if there is one, lies in a single row, so the other row is all $b$. No row is all $1-b$, because that would need two entries that differ from $b$. So the \dt has a certificate that forces the label $b$ and none that forces $1-b$. Hence $h_{v_{b,e}}=h_b$, and $\H_\V=\H$.

Next, we give the strategy for the realizable case. Fix $v^*\in\V$ and a sequence of inputs, and let $h^*=h_{v^*}=h_b$. In each round, both debaters choose their certificates uniformly and independently, and the \dt proposes the majority of the earlier feedback $\vy_1,\ldots,\vy_{t-1}$, breaking ties arbitrarily. Since each table of $v^*$ has at most one entry equal to $1-b$, the values $\vy_t$ are independent, and each equals $b$ with probability at least $3/4$. After $s$ rounds, the majority is wrong only if at most $s/2$ of the $s$ values equal $b$, while their expected number is at least $3s/4$. By Hoeffding's inequality, this has probability at most $e^{-2(s/4)^2/s}=e^{-s/8}$. Summing over rounds gives $\ex[\Mi(T)]\le\sum_{s\ge0}e^{-s/8}=1/(1-e^{-1/8})<9$. The randomization is needed here. If the debaters' choices were deterministic, the environment could put the exceptional entry of each new table where they look.

Finally, we prove the lower bound in the unrealizable case. The idea is that a random entry of a table in $\V$ shows the true label $b$ with probability $3/4$, while a misspecified round can always show $1-b$. If each round is misspecified with probability $1/3$, the two effects cancel, and every feedback value is a fair coin. Since every input is new, the debaters cannot tell which rounds are misspecified, so they learn nothing about $b$.

Let $x_t=t$ and $m=3k$, and consider the following random instance, drawn before the interaction. Sample $b$ uniformly and set $h^*=h_b$. Independently for each $t\le m$, with probability $1/3$ make round $t$ misspecified and set every entry of $v^*(x_t,\cdot,\cdot)$ to $b$, so that $\vy_t=1-b$ whatever the certificates. Otherwise, set every entry of $v^*(x_t,\cdot,\cdot)$ to $b$ except one uniformly random entry $e_t\in\C\times\C$, which we set to $1-b$. No later round is misspecified, and for every input $x>m$, every entry of $v^*(x,\cdot,\cdot)$ is $b$. Every table of $v^*$ has at most one entry that differs from $b$, so $v^*\in\V$ validates $h^*$. Let $K$ be the number of misspecified rounds, so $K\sim\mathrm{Binomial}(m,1/3)$. In particular, $\ex[K]=m/3=k$.

Fix a strategy (potentially randomized) of the debaters, which may depend on $k$ and $T$. We show that $\ex[\Mi(m)]=m/2$. Let $I_{t-1}$ consist of the debaters' random bits and the transcript of the first $t-1$ rounds. On a misspecified round $\vy_t=1-b$, and on any other round $\vy_t=b$ unless $(\ct_t,\cf_t)=e_t$. Hence, for each $a\in\{0,1\}$ and any certificates, $\Pr(\vy_t=a\mid h^*=h_a,I_{t-1})=\frac23\cdot\frac34+\frac13\cdot0=\frac12$. So whichever of $h_0$ and $h_1$ is the true hypothesis, each $\vy_t$ is a fair coin independent of $I_{t-1}$, which means $\Pr(\hy_t\ne\ty_t)=1/2$.  Summing over $t\le m$ gives $\ex[\Mi(m)]=m/2$. Since $m=3k\le T$, we get $\ex[\Mi(T)]\ge\ex[\Mi(m)]=\frac32k$, while the expected number of misspecified rounds is $\ex[K]=k$.
\end{proof}
